\section{Янг–Миллс и Ходж}\label{sec:ymhodge}

Две задачи Клэя~\cite{clay-ym,clay-hodge} на языке этой статьи читаются как запреты вечного
двигателя. В Янг--Миллсе запрещённый объект --- безмассовый спектр: бесконечная башня всё более
дешёвых возбуждений над вакуумом. В Ходже --- неоплаченный заряд, бесконечно убывающий по натуральной
высоте. В обеих ветвях мы доказываем структурную половину, не зависящую от конкретного аналитического
объекта, которого в mathlib нет, и называем красный вход-якорь, единственно и остающийся. Lean-источники
--- \code{Engine/YangMillsFront.lean} и \code{Engine/HodgeFront.lean}.

\subsection{Янг–Миллс: безмассовость как вечный двигатель}

Задача о массовой щели просит доказать, что у квантовой калибровочной теории есть \emph{щель} ---
минимальная положительная цена возбуждения над вакуумом. Её отрицание --- \emph{безмассовость}, и
прочтение ветви состоит в том, что безмассовость есть вечный двигатель: работа, черпаемая из вакуума
без нижнего предела. Подчёркиваем: задачу Клэя мы \emph{не} решаем и решённой не объявляем. В mathlib
нет калибровочной теории --- ни гамильтониана, ни настоящего спектра. Мы работаем с абстрактной
\emph{спектральной моделью} \code{SpectralModel} (множество уровней энергии
$\mathrm{Energy}\subseteq\R$ с вакуумом $0$ и неотрицательностью) и доказываем не зависящую от теории
половину.

На ней \emph{масс-щель} $\mathrm{MassGap}(S)$ и \emph{безмассовость} $\mathrm{Gapless}(S)$ ---
точные отрицания друг друга.

\begin{theorem}[\code{not\_massGap\_iff\_gapless}]\label{thm:not_massGap_iff_gapless}
\green{}\ Отсутствие щели эквивалентно безмассовости, и без всякого выбора:
\[
  \neg\,\mathrm{MassGap}(S)\iff\mathrm{Gapless}(S),
\]
где $\mathrm{MassGap}(S):=\exists\,\Delta>0,\ \forall E\in S.\mathrm{Energy},\ E\neq 0\Rightarrow\Delta\le E$
и $\mathrm{Gapless}(S):=\forall\varepsilon>0,\ \exists E\in S.\mathrm{Energy},\ 0<E<\varepsilon$.
\end{theorem}

Двигателем безмассовость делает то, что из неё строится \emph{лестница} --- halving-цепочка
положительных уровней, на каждом шаге падающая не меньше чем вдвое, $2\cdot E(n+1)\le E(n)$. Это ровно
$\R$-контрпример положительной работы без нижнего предела, бесконечный мультипликативный спуск в
вещественных числах, перед которым ацикличность бессильна.

\begin{theorem}[\code{gapless\_iff\_nonempty\_ladder}]\label{thm:gapless_iff_nonempty_ladder}
\green{}\ Безмассовость эквивалентна существованию лестницы: $\mathrm{Gapless}(S)$ выполнено тогда и
только тогда, когда тип $\mathrm{GaplessLadder}(S)$ лестниц непуст, где лестница --- это
последовательность $E:\N\to\R$ с $E(n)\in S.\mathrm{Energy}$, $E(n)>0$ и $2\cdot E(n+1)\le E(n)$ для
всех $n$.
\end{theorem}

\emph{Идея доказательства.} Прямо лестница строится выбором, зеркаля извлечение нуля в римановой
ветви; обратно лестница опровергает любой кандидат $\Delta$ геометрической мажорантой. Пока лестница
живёт в $\R$, ацикличность её не касается.

Двигатель убивается ровно тогда, когда лестницу удаётся \emph{отразить в натуральные числа}. Это
физический вход ветви, и он по-настоящему открыт.

\begin{definition}[\code{QuantizationLaw}]\label{def:QuantizationLaw}
\red{}\ Красный по-модельный вход (паттерн \code{EnergyBalanceLaw}): каждое положительное состояние
несёт натуральный ранг, а мультипликативный спуск энергии отражается в \emph{строгий} спуск ранга,
\[
  \mathrm{QuantizationLaw}(S):=\exists\,\mathrm{rank}:\mathrm{PositiveState}(S)\to\N,\
    \forall x,y,\ 2\cdot(y:\R)\le(x:\R)\Rightarrow\mathrm{rank}(y)<\mathrm{rank}(x).
\]
\end{definition}

При законе лестница из вещественной становится натуральной и обрывается.

\begin{theorem}[\code{no\_quantizedLadder}]\label{thm:no_quantizedLadder}
\green{}\ Квантованной лестницы не существует:
$\mathrm{QuantizationLaw}(S)\wedge\mathrm{GaplessLadder}(S)\Rightarrow\bot$.
\end{theorem}

\emph{Идея доказательства.} Квантованная лестница была бы бесконечной строго убывающей цепью
натуральных рангов --- вечным двигателем Евклида --- и убивается корнем всего репозитория,
\code{no\_infinite\_descent} (\cref{sec:engine}). Убийца здесь --- чистая фундированность $\N$; никакой
фальшивой гипотезы «двигателей нет» не подмешано.

\begin{theorem}[\code{massGap\_of\_quantizationLaw}]\label{thm:massGap_of_quantizationLaw}
\green{}\ Квантование влечёт масс-щель: $\mathrm{QuantizationLaw}(S)\Rightarrow\mathrm{MassGap}(S)$.
\end{theorem}

\emph{Идея доказательства.} Нет щели --- есть безмассовость (\cref{thm:not_massGap_iff_gapless}),
безмассовость даёт лестницу (\cref{thm:gapless_iff_nonempty_ladder}), квантование превращает её в
натуральный спуск, и \cref{thm:no_quantizedLadder} замыкает противоречие. Это зелёная условная теорема
ветви: спектральное квантование $\Rightarrow$ масс-щель.

Как и в P/NP, возникает соблазн добавить четвёртую границу первопричины; машина этого не позволяет, и
причина здесь глубже.

\begin{theorem}[\code{quantizationLaw\_iff\_massGap}]\label{thm:quantizationLaw_iff_massGap}
\green{}\ Для каждой модели закон квантования эквивалентен наличию щели, зелено и без всякой границы:
$\mathrm{QuantizationLaw}(S)\iff\mathrm{MassGap}(S)$.
\end{theorem}

Это форма \emph{осуждённого} моста, а не честного риманова зеркала: там эквивалентность «закон $\iff$
цель» держалась лишь \emph{под} границей, здесь она зелена сама по себе. Декретировать закон квантования
означало бы переименовать цель, а не добавить причину. Трилемма подтверждает это со всех сторон:
универсальная форма опровержима --- подделанная безмассовая модель $\{0\}\cup\{2^{-n}\}$ несёт явную
лестницу \code{cookedLadder}, и декрет взорвал бы карантин
(\code{ymManifestationLaw\_refutes\_boundary}, \green{}); экзистенциальная форма уже доказана
подделанной моделью со щелью $\{0,1\}$, так что декрет был бы пуст. Недостаёт не утверждения, а
\red{} якоря данных --- построенного спектра или гамильтониана настоящей неабелевой теории, которого в
mathlib нет.

Единственное \yellow{} содержание ветви --- то, что декрет утверждает на своём масштабе. При $A\le 4$
запас отклонений реален (\code{smallScale\_deviationSupply}, \green{}); на декретированном $A\ge 5$
запаса нет ни при каком пороге (\code{decreedScale\_no\_deviationSupply}, \yellow{}, условно на
\axiom{}): мир первопричины имеет щель на языке запасов. Умышленный детектор взрыва,
\code{quarantine\_inconsistent\_if\_supply\_at\_every\_scale} (\yellow{}), держит точку взрыва на виду.

Эпистемическое дополнение (\code{Engine/YangMillsEpistemic.lean}) закрывает внутреннюю дверь, как в
Коллатце и P/NP: самообоснование закона квантования саморазрушается
(\code{no\_internalisedYMGround}, \green{}), откуда \code{ymCause\_unknowable} (\green{}). За
противоречие платит настоящая конструкция: из отсутствия щели строится halving-лестница, несущая
настоящий вечный двигатель на $\R$ (\code{not\_massGap\_carries\_real\_engine}, \green{}), законный на
$\R$; стена возникает лишь когда ранг переводит его в $\N$-спуск, сгорающий о
\code{no\_perpetual\_engine\_on\_nat}. Всё собрано в \code{ym\_locked\_behind\_engine\_status}
(\green{}); таинт репозитория не растёт.

\begin{remark}
Это \emph{не} решение задачи Клэя и \emph{не} Гёдель: ни неполноты, ни неподвижной точки, только
фундированность, и вся эпистемика --- внутримодельная. О настоящей квантовой теории Янга--Миллса
теоремы не говорят ничего; она остаётся \red{}. Условность названа: закон квантования --- \red{}
по-модельный вход, и всё «под законом» здесь импликация; тавтологизация через
\cref{thm:quantizationLaw_iff_massGap} не скрыта --- именно поэтому единственный выход внешний.
\end{remark}

\subsection{Ходж: квантованные заряды и оплата циклами}

Класс Ходжа мы читаем как \emph{квантованный заряд} --- рациональный класс когомологий выровненного
типа $(p,p)$ --- а алгебраический цикл как его \emph{оплату}. Гипотеза Ходжа тогда --- полнота: всякий
квантованный выровненный заряд оплачен; неоплаченный заряд открыл бы бесконечный строгий спуск по
натуральной высоте, запрещённый вечный двигатель. Снова честная граница: теории Ходжа в mathlib
\emph{нет} --- ни классов Ходжа, ни алгебраических циклов, ни когомологий многообразий. Мы работаем с
абстрактным леджером \code{HodgeLedger} (классы, предикаты \code{IsHodge} и \code{IsAlgebraic}, высота
$\mathrm{height}:\mathrm{Cls}\to\N$ и якорь квантования «оплачен $\iff$ высота $0$»). Задачу мы не
решаем; остаток --- именованный вход \code{DescentLaw}.

Заголовок ветви в том, что здесь двигатель мёртв \emph{безусловно}.

\begin{theorem}[\code{isEmpty\_unpaidDescentChain}]\label{thm:isEmpty_unpaidDescentChain}
\green{}\ Для каждого Ходж-леджера $S$ тип $\mathrm{UnpaidDescentChain}(S)$ пуст: бесконечной строго
убывающей по высоте цепи неоплаченных зарядов не существует ни в одной модели. Список аксиом --- лишь
\code{[propext, Quot.sound]}, даже без выбора: чистый спуск при $A=1$.
\end{theorem}

В отличие от Янг--Миллса, где вещественная лестница жила при безмассовости и убивалась
\emph{декретируемым} законом квантования, здесь квантование встроено в модель --- высота есть
натуральное число --- поэтому двигатель мёртв в каждой модели. Оспаривается не то, квантован ли заряд
(да, по построению), а существование \emph{шагов} оплаты.

\begin{theorem}[\code{hodgeProperty\_of\_descentLaw}]\label{thm:hodgeProperty_of_descentLaw}
\green{}\ Если $\mathrm{DescentLaw}(S)$, то $\mathrm{HodgeProperty}(S)$: закон спуска влечёт гипотезу
Ходжа модели.
\end{theorem}

\emph{Идея доказательства.} Сильная индукция по высоте, выбор не нужен; то же следует цепным путём
через \code{descentSeq}. Это зелёная условная теорема ветви: закон спуска --- оплата циклами.

\begin{theorem}[\code{no\_unpaid\_lawful\_model}]\label{thm:no_unpaid_lawful_model}
\green{}\ Нет леджера $S$, удовлетворяющего одновременно $\mathrm{DescentLaw}(S)$ и
$\mathrm{Nonempty}(\mathrm{UnpaidClass}(S))$: модели «закон плюс неоплаченный заряд» не существует.
\end{theorem}

Обратное честно вакуумно: \code{descentLaw\_of\_hodgeProperty} доказано пустым квантором (если всё
оплачено, неоплаченных зарядов нет), в отличие от Янг--Миллса, где обратный ранг строился по-настоящему.
Отсюда --- решающий аудит.

\begin{theorem}[\code{descentLaw\_iff\_hodgeProperty}]\label{thm:descentLaw_iff_hodgeProperty}
\green{}\ Для каждого леджера $S$: $\mathrm{DescentLaw}(S)\iff\mathrm{HodgeProperty}(S)$, зелено и без
всякой границы.
\end{theorem}

Снова мост осуждённого: декретировать закон означало бы декретировать гипотезу дословно. Шестого поля
декрета быть не может; недостаёт \red{} якоря данных --- настоящих рациональных $(p,p)$-классов.
Трилемма это подтверждает, и одна её грань здесь уникальна. Перенос формы манифестации из Янг--Миллса
над \emph{цепями} вырождается в зелёную теорему.

\begin{theorem}[\code{hodgeChainManifestationLaw\_green}]\label{thm:hodgeChainManifestationLaw_green}
\green{}\ $\mathrm{HodgeChainManifestationLaw}$: для каждого леджера $S$ и каждой цепи
$C:\mathrm{UnpaidDescentChain}(S)$ имеем $\mathrm{ChainManifests}(C)$. Поскольку по
\cref{thm:isEmpty_unpaidDescentChain} такой цепи не существует, закон над цепями вакуумно истинен.
\end{theorem}

Это вакуумность типа «пустой свидетель», а не несовместимость: из пустого типа свидетелей нельзя
собрать третью грань трилеммы, которой нужен свидетель, \emph{предъявленный} в зелёном. Честный свидетель
здесь не цепь, а единственный неоплаченный класс, \code{cookedUnpaidClass}, предъявленный по построению
без выбора.

\begin{theorem}[\code{hodgeManifestationLaw\_refutes\_boundary}]\label{thm:hodgeManifestationLaw_refutes_boundary}
\green{}\ Если выполнено $\mathrm{HodgeManifestationLaw}$, то
$\neg\,\mathrm{TheStrictLastStep00Obligation}$: форма манифестации закона несовместима с принятой
границей карантина.
\end{theorem}

Единственное \yellow{} ветви --- растяжка,
\code{quarantine\_inconsistent\_if\_hodgeManifestationLaw\_decreed} (условно на \axiom{}): декрет формы
манифестации взорвал бы карантин ровно здесь, и её безопасность проверена.

Эпистемическое дополнение (\code{Engine/HodgeEpistemic.lean}) закрывает внутреннюю дверь.
Самообоснование по-модельного закона спуска саморазрушается (\code{no\_internalisedHodgeGround}).

\begin{theorem}[\code{hodgeCause\_unknowable}]\label{thm:hodgeCause_unknowable}
\green{}\ Для каждого леджера $S$: $\neg\,\mathrm{InternalKnowledgeOfHodgeCause}(S)$ --- внутреннее
самообоснование закона спуска невозможно в каждой модели, безусловно и без всяких аксиом.
\end{theorem}

У противоречия есть второй, двигательный путь: под законом неоплаченный заряд разворачивается в
настоящую бесконечную цепь спуска (\code{unpaidDescentChain\_of\_descentLaw}) и проектируется высотой в
вечный двигатель на $\N$ (\code{unpaidClass\_under\_law\_carries\_perpetual\_engine}), построенный
по-настоящему, а не \emph{ex falso} (\code{internalisedHodgeGround\_builds\_engine}) и сжигаемый о
\code{no\_perpetual\_engine\_on\_nat}. Универсальная форма уже опровергнута подделкой
(\code{hodgeUniversal\_forged\_refutation}), так что единственный живой вариант --- по-модельный вход,
который по \cref{thm:descentLaw_iff_hodgeProperty} равен цели дословно.

\begin{theorem}[\code{hodge\_locked\_behind\_engine\_status}]\label{thm:hodge_locked_behind_engine_status}
\green{}\ Для каждого леджера $S$:
\begin{equation}\label{eq:hodge_locked}
  (\neg\,\mathrm{HodgeDescentLawUniversal})\ \wedge\ \mathrm{IsEmpty}(\mathrm{UnpaidDescentChain}(S))
  \ \wedge\ (\neg\,\mathrm{InternalKnowledgeOfHodgeCause}(S))
  \ \wedge\ \bigl(\mathrm{DescentLaw}(S)\iff\mathrm{HodgeProperty}(S)\bigr).
\end{equation}
Универсальная опровергнута, цепи пусты в каждой модели, внутреннее знание невозможно, по-модельный
закон эквивалентен гипотезе модели --- целиком зелено; таинт репозитория не растёт.
\end{theorem}

\begin{remark}
Это \emph{не} решение гипотезы Ходжа и \emph{не} Гёдель: ни неполноты, ни неподвижной точки, только
внутримодельная эпистемика и $\N$-спуск. Самообосновывающая пара формально --- «обоснование $+$ его
отрицание»; её содержание оплачено зелёным отождествлением обеих сторон с содержанием (закон $\iff$
гипотеза модели по \cref{eq:hodge_locked}; отрицание $\iff$ предъявимость неоплаченного заряда,
\code{not\_hodgeProperty\_iff\_unpaidClass}) и настоящим двигательным путём. Цвет Ходжа на общем
правиле: когда квантование встроено в модель, двигатель мёртв безусловно и гипотеза сжимается к
единственному вопросу --- существуют ли шаги оплаты.
\end{remark}
